\clearpage
\begingroup
\setlength{\parskip}{1pt}
\titlespacing*{\section}{0pt}{0.8em}{0.45em}
\titlespacing*{\subsection}{0pt}{0.45em}{0.25em}
\section*{Reading an Instruction Description}\label{page:reading-an-instruction-description}
\phantomsection
\addcontentsline{toc}{section}{Reading an Instruction Description}
\subsection*{Entry Organization}
Each instruction entry presents its Operation, then Assembler Syntax, Attributes, and structured status
metadata. Detailed Semantics follows when more space is needed, before the concrete encodings. Brief
summaries appear only in the instruction-set summary tables. Each mnemonic begins on a new page so its
encoding diagrams, field explanations, and side-effect text stay together as one reference unit.

\subsection*{Mnemonics and Suffixes}
The B, W, L, and Q suffixes select 8-, 16-, 32-, and 64-bit integer sizes; S and D select 32- and 64-bit
floating-point scalar sizes where defined.

When an instruction encoding includes a size selector z, its definition maps z to the suffix set shown for that encoding. A
mnemonic without a size suffix has a single architecturally defined size or no data-size operand.
Conditional mnemonic variants use the condition names listed in the Flags and Condition Codes chapter.

CMPJcc, DJcc, IJcc, Jcc, MOVcc, SETcc, TESTJcc, and FMOVcc place the condition suffix in the mnemonic. An encoding with
only one architectural size omits the size suffix.

\subsection*{Operand and Status Notation}
Rn names an Rn register selected from the general-purpose registers R0 through R15; SP and PC are special registers
outside that encoding namespace. An
\texttt{<ea>} operand names a compact effective-address field and any appended EA payload.

FLAGS, STATUS, FFLAGS, and FSTATUS are architectural state registers. Instruction descriptions state whether an
encoding updates, preserves, reads, or ignores each status field.

In instruction pseudocode, Operand Read applies the addressing mode and operation size. Operand Write stores the
selected-size result; a sub-Q Rn write preserves upper bits unless the instruction defines another rule. EA Address
computes an address without reading memory. A control-target memory operand reads its declared-width target, while a
control-target register or immediate operand supplies the target value directly. A selector may name an Rn register, a
segment register, or an encoded immediate. FLAGS ZNCV refers to the integer Z, N, C, and V bits, while FFLAGS names
the accrued floating-point exception flags.

\subsection*{Encodings}
In an encoding diagram, fixed 0 and 1 fields identify framing, opcode-selection, or reserved values, and the notes
below decode lettered fields. An <ea> field is the compact effective-address field and may select one or more
independently determined EA payloads in operand order.

In each instruction entry, Encodings presents the encoding variants, operand and
effective-address grammar, and size selectors.

For an extended encoding, the four-bit L field selects a 3+L-byte encoded instruction length. The required
instruction length is the opcode-space length plus the selected operand payload lengths. Every larger encoded
instruction length through 18 bytes is valid, and the trailing bytes are uninterpreted padding. The
\texttt{LEN n,} spelling records an explicit encoded instruction length.
\endgroup
\clearpage
